\begin{thebibliography}{99}
\addcontentsline{toc}{chapter}{Bibliographie}

\bibitem{coco} T.-Y. Lin, M. Maire, S. Belongie et al., \emph{Microsoft COCO: Common Objects in Context}, European Conference on Computer Vision (ECCV), 2014. \url{https://cocodataset.org}

\bibitem{hog} N. Dalal et B. Triggs, \emph{Histograms of Oriented Gradients for Human Detection}, IEEE Conference on Computer Vision and Pattern Recognition (CVPR), 2005.

\bibitem{svm} C. Cortes et V. Vapnik, \emph{Support-Vector Networks}, Machine Learning, vol.~20, 1995.

\bibitem{rf} L. Breiman, \emph{Random Forests}, Machine Learning, vol.~45, 2001.

\bibitem{knn} T. Cover et P. Hart, \emph{Nearest Neighbor Pattern Classification}, IEEE Transactions on Information Theory, vol.~13, 1967.

\bibitem{xgboost} T. Chen et C. Guestrin, \emph{XGBoost: A Scalable Tree Boosting System}, ACM SIGKDD, 2016.

\bibitem{rcnn} R. Girshick, J. Donahue, T. Darrell et J. Malik, \emph{Rich Feature Hierarchies for Accurate Object Detection and Semantic Segmentation}, CVPR, 2014.

\bibitem{dijkstra} E.~W. Dijkstra, \emph{A Note on Two Problems in Connexion with Graphs}, Numerische Mathematik, vol.~1, 1959.

\bibitem{osm} Contributeurs OpenStreetMap, données sous licence ODbL. \url{https://www.openstreetmap.org}

\bibitem{sklearn} F. Pedregosa et al., \emph{Scikit-learn: Machine Learning in Python}, Journal of Machine Learning Research, vol.~12, 2011.

\bibitem{pytorch} A. Paszke et al., \emph{PyTorch: An Imperative Style, High-Performance Deep Learning Library}, NeurIPS, 2019.
\end{thebibliography}

\appendix

\chapter{Sources des illustrations}
\label{annexe:sources}

Les figures produites par nos notebooks et notre simulateur sont les nôtres. Les autres illustrations viennent de sources externes et servent seulement à illustrer le propos.

\begin{center}
\small
\begin{tabularx}{\textwidth}{>{\raggedright\arraybackslash}p{5.6cm}Y}
  \toprule
  \textbf{Illustration} & \textbf{Source} \\
  \midrule
  Scénarios 1 à 4 & \acompleter{Source et licence à indiquer} \\
  Environnements virtuels (figure~\ref{fig:inspiration}) & \acompleter{Source et licence à indiquer} \\
  Détection par R-CNN (figure~\ref{fig:detection}) & \acompleter{Source et licence à indiquer} \\
  Modèles 3D des classes & Captures de sites de modèles 3D, dont fab.com (moto) et turbosquid.com (voiture) ; \acompleter{auteurs à indiquer} \\
  Flux de développement, répartition des efforts & Schémas réalisés pour le projet \\
  Statistiques, recadrages et comparaison des modèles & Nos notebooks, à partir de COCO 2017 \\
  Capture du simulateur & Notre simulateur ; carte © contributeurs OpenStreetMap (ODbL) \\
  \bottomrule
\end{tabularx}
\end{center}

\chapter{Usage de l'intelligence artificielle}

Le simulateur et la structure des notebooks ont été préparés avec l'aide de Claude Code, et la rédaction de ce dossier a aussi été assistée par cet outil. Les résultats préliminaires de la section~\ref{sec:preuves} et les chiffres du chapitre~\ref{chap:donnees} viennent de notebooks complétés et exécutés par Claude Code sur la branche \texttt{correction} du dépôt, à la demande de l'équipe. Nous les présentons comme des preuves de faisabilité. Leur place dans le travail noté reste à définir avec notre encadrant.
